
      \documentclass[fleqn]{article}      
      \usepackage{amsmath}
      \usepackage{fontspec}
      \usepackage{kotex} % 한국어 지원  

      \begin{document}      
      쉬운 객관식 문제: \\\\  
1. 함수의 정의는 무엇인가요? \\\\  
   a) 두 집합 사이의 관계 \\\\  
   b) 수의 곱셈 \\\\  
   c) 도형의 넓이 \\\\  

2. 다음 중 함수의 그래프가 아닌 것은? \\\\  
   a) 직선 \\\\  
   b) 포물선 \\\\  
   c) 원 \\\\  

3. 주어진 집합 $A = \{1, 2, 3\}$에서 함수 $f : A \to A$의 값으로 가능한 조합의 수는? \\\\  
   a) 6 \\\\  
   b) 9 \\\\  
   c) 3 \\\\  

중간 객관식 문제: \\\\  
1. 함수 $f(x) = 2x + 3$의 기울기는 얼마인가요? \\\\  
   a) 2 \\\\  
   b) 3 \\\\  
   c) 5 \\\\  

2. 함수 $f(x) = x^2 - 4$의 정의역이 $[-2, 2]$일 때, $f(1)$의 값은? \\\\  
   a) 0 \\\\  
   b) -3 \\\\  
   c) -1 \\\\  

3. 함수 $g(x) = \frac{1}{x}$의 정의역은 무엇인가요? \\\\  
   a) 모든 실수 \\\\  
   b) $x \neq 0$ \\\\  
   c) $x > 0$ \\\\  

어려운 객관식 문제: \\\\  
1. 함수 $f(x) = x^3 - 3x$의 극대값과 극소값을 구하시오. \\\\  
   a) 극대: 2, 극소: -2 \\\\  
   b) 극대: 0, 극소: -3 \\\\  
   c) 극대: 3, 극소: 1 \\\\  

2. 함수 $h(x) = 2^x$의 그래프가 지나는 점은? \\\\  
   a) $(0, 1)$ \\\\  
   b) $(1, 2)$ \\\\  
   c) $(2, 4)$ \\\\  

3. $f(x) = 3x + 2$와 $g(x) = -x + 1$의 교점을 구하시오. \\\\  
   a) $(1, 5)$ \\\\  
   b) $(0, 2)$ \\\\  
   c) $(2, 8)$ \\\\  

쉬운 주관식 문제: \\\\  
1. 함수의 정의를 설명하시오. \\\\  
2. $y = 3x + 2$에서 기울기의 의미를 설명하시오. \\\\  
3. 함수의 그래프를 그릴 때 필요한 정보는 무엇인가요? \\\\  
4. 주어진 집합 $A = \{1, 2, 3\}$에서 가능한 함수의 수를 구하시오. \\\\  
5. 함수 $f(x) = x + 1$의 값을 $x = 3$일 때 구하시오. \\\\  

중간 주관식 문제: \\\\  
1. 함수 $f(x) = x^2 + 3x + 2$의 최소값을 구하시오. \\\\  
2. 함수 $g(x) = \sqrt{x}$의 정의역을 구하시오. \\\\  
3. 두 함수 $f(x) = 2x + 1$과 $g(x) = 3x - 4$의 교점을 구하시오. \\\\  
4. 함수 $h(x) = -x^2 + 4$의 최대값을 구하시오. \\\\  
5. 다음 함수의 주기를 구하시오: $f(x) = \sin(x)$. \\\\  

어려운 주관식 문제: \\\\  
1. 함수 $f(x) = x^3 - 6x^2 + 9x$의 극대값과 극소값을 구하시오. \\\\  
2. $f(x) = e^x$의 미분을 구하시오. \\\\  
3. 두 함수가 서로 역함수가 되기 위한 조건을 설명하시오. \\\\  
4. 함수 $f(x) = \log(x)$의 정의역을 설명하시오. \\\\  
5. $f(x) = 2x + 3$와 $g(x) = -x + 1$의 합성함수 $f(g(x))$를 구하시오. \\\\  

      
      \end{document} 
  